\scanpage{30}
\mspage{4}
\unsure{en} la normalisant\inserted{p. ex.} de telle façon que la longueur du cercle $|z|=1$ soit \deleted{$2\pi$} $1$, on la définit sur toutes les $X\simeq\mathbb C^*$ en $=$ fait. D'ailleurs l'anti-automorphisme $z\longmapsto\bar z$ \deleted{\unsure{qui}} conserve la mesure, donc si $X\simeq\mathbb C^*$, \uline{tt} autom. conforme respecte la mesure canonique.

\margin{C'est la métrique déduite de \unsure{C! définie|celle définie} par hermitienne canonique $z\bar z$ sur $\mathbb C$, via $z\longmapsto\exp 2i\pi z$, de $\mathbb C$ sur $\mathbb C^*$.}

Il en résulte que dans le cas non orienté $X\simeq$ plan projectif réel moins $1$ pt, on peut descendre à partir de $X^{or}$ et on trouve aussi une métrique canonique plate, stable par tous les automorphismes conformes.

$2''$) Cas $X\simeq\mathbb E^1_{\mathbb C}$. Les métriques riemanniennes sur une \deleted{\unsure{courbe}} droite vectorielle complexe $L$, associées aux structures hermitiennes\inserted{déf. pos.} sur $L$, (elles sont toutes proportionnelles) ne sont pas invariantes par autom \deleted{\unsure{de}}\inserted{de} $L$, mais seulement par le sous-groupe engendré par les translations, et les \deleted{\unsure{autom.}} homothéties par des $u\in\mathbb U=\{z\in\mathbb C\mid |z|=1\}$; elles sont quasi-invariantes par autom. quelconques (lesquels sont \uline{affines}). Donc sur une $X\simeq\mathbb E^1_{\mathbb C}$, on peut définir, non \uline{une} métrique
% archive foliation: 32
